\chapter{Paradigmi emergenti}\label{ch-11}

\lecturemeta{Capitolo originale: 11 --- Emerging Paradigms in Computing \quad\textbullet\quad Antonio Brogi}{Appunti di Emiliano Sescu (github.com/Faxatos), cap. 11}

\section{Cloud-Edge Continuum}\label{cloud-edge-continuum}

\begin{boxdefinition}{Cloud-Edge Continuum}

Il \textbf{Cloud-Edge Continuum} unisce i punti di forza
dell'\textbf{Edge Computing} (elaborare i dati vicino a dove vengono
prodotti) e del \textbf{Cloud Computing}, estendendo i servizi cloud
fino all'\textbf{Internet of Things} (IoT). Il risultato è
un'infrastruttura \textbf{distribuita ed eterogenea}.

\end{boxdefinition}

Vantaggi principali:

\begin{itemize}
\tightlist
\item
  \textbf{Potenza di calcolo}: si possono usare sia le risorse del cloud
  sia i dispositivi edge.
\item
  \textbf{Connettività}: comunicazione migliore tra dispositivi, con
  trasferimento di dati più fluido.
\item
  \textbf{Bassa latenza}: elaborare i dati vicino alla loro fonte riduce
  i ritardi e rende il sistema più reattivo.
\end{itemize}

Le applicazioni del futuro saranno soprattutto \textbf{a microservizi in
container}, eseguite su un'infrastruttura continua Cloud-Edge. Gestirle
però è difficile, perché sia l'infrastruttura sia le applicazioni
cambiano di continuo:

\begin{itemize}
\tightlist
\item
  \textbf{Cambiamenti dell'infrastruttura}: il carico dei nodi può
  spostarsi, latenza e banda possono variare, i nodi possono entrare o
  uscire dalla rete all'improvviso, e possono esserci interruzioni
  temporanee di connessione.
\item
  \textbf{Cambiamenti delle applicazioni}: il codice e i requisiti
  cambiano e richiedono adattamenti rapidi.
\end{itemize}

Per questo il deploy delle applicazioni va \textbf{gestito di continuo},
anche dopo il primo avvio.

\subsection{Monitoraggio}\label{monitoraggio}

Serve un monitoraggio efficace sia delle applicazioni sia
dell'infrastruttura: un sistema \textbf{leggero}, \textbf{tollerante ai
guasti} e capace di adattarsi ai cambiamenti.

Un metodo efficace è il \textbf{continuous reasoning}: si analizzano
sistemi grandi concentrandosi sui \textbf{cambiamenti recenti} e
\textbf{riusando i risultati precedenti} dove possibile. Nel continuum
bisogna considerare insieme i cambiamenti dell'applicazione e
dell'infrastruttura per decidere se sostituire, spostare, riavviare o
scalare i servizi.

Un esempio recente è \textbf{FogBrainX}
(\href{https://github.com/di-unipi-socc/fogbrainx}{GitHub}), che
confronta le differenze nelle specifiche dell'applicazione e nei dati
monitorati dell'infrastruttura. Aiuta a decidere dove collocare i
servizi (\emph{placement}):

\begin{itemize}
\tightlist
\item
  adattandosi ai cambiamenti dell'infrastruttura (es. risorse dei nodi o
  qualità della rete), che possono richiedere di spostare dei servizi;
\item
  adattandosi ai cambiamenti nei requisiti dei servizi (software,
  hardware, IoT) o nelle esigenze di comunicazione, che possono
  richiedere aggiornamenti;
\item
  gestendo l'aggiunta o la rimozione di servizi o di requisiti di
  comunicazione nelle specifiche dell'applicazione.
\end{itemize}

Prese le decisioni, FogBrainX le passa a \textbf{FogArm}, che esegue i
comandi di gestione sull'infrastruttura Cloud-IoT.

\subsection{Gestione decentralizzata}\label{gestione-decentralizzata}

Per la gestione decentralizzata ci sono due approcci interessanti:

\begin{itemize}
\tightlist
\item
  \textbf{Gestione osmotica} (\emph{osmotic management}): i servizi
  dell'applicazione si adattano in base alle risorse disponibili e alle
  esigenze dell'applicazione. Le politiche di gestione possono
  disinstallare, spostare o scalare le applicazioni in tempo reale.
\item
  \textbf{Gestione decentralizzata} ispirata al comportamento dei
  \textbf{batteri}:

  \begin{itemize}
  \tightlist
  \item
    ogni istanza dell'applicazione ha un proprio \textbf{agente di
    gestione} (una specie di mini-unità di gestione);
  \item
    \textbf{regole semplici} fanno scattare azioni (es. disinstallare o
    replicare) in base ai dati monitorati;
  \item
    dai comportamenti locali \textbf{emerge} un comportamento globale
    flessibile, utile per infrastrutture grandi, anche se più complesso
    da gestire.
  \end{itemize}
\end{itemize}

\section{Quantum Software
Engineering}\label{quantum-software-engineering}

\begin{boxdefinition}{Quantum Software Engineering (QSE)}

La \textbf{Quantum Software Engineering} applica solidi principi di
ingegneria allo sviluppo, all'esercizio e alla manutenzione del
\textbf{software quantistico} e della sua documentazione. L'obiettivo è
creare software quantistico affidabile, efficiente sui computer
quantistici e con costi contenuti.

\end{boxdefinition}

Il \textbf{Manifesto di Talavera} per la QSE indica alcuni principi
importanti, tra cui:

\begin{itemize}
\tightlist
\item
  la QSE deve essere compatibile con diversi linguaggi e tecnologie di
  programmazione quantistica;
\item
  bisogna integrare calcolo \textbf{classico e quantistico}, usando i
  computer quantistici soprattutto per i compiti specifici in cui sono
  migliori (es. i problemi di fattorizzazione).
\end{itemize}

\subsection{Quantum Broker}\label{quantum-broker}

Un esempio pratico di QSE è il \textbf{Quantum Broker}, che risponde
alla domanda: ``\textbf{Quale computer quantistico devo usare per
eseguire il mio algoritmo?}''. È utile soprattutto per chi non conosce
bene i vari provider quantistici.

Anche dopo aver scelto un provider e avviato gli algoritmi, possono
esserci problemi:

\begin{itemize}
\tightlist
\item
  \textbf{Disponibilità}: cosa succede se il computer quantistico
  diventa non disponibile durante l'esecuzione?
\item
  \textbf{Requisiti}: come bilanciare costo, tempo e accuratezza?
\item
  \textbf{Personalizzazione}: come adattare il modo in cui si sceglie il
  computer quantistico alle proprie esigenze?
\end{itemize}

\begin{boxdefinition}{Shot e distribuzione degli shot}

Un circuito quantistico ha un risultato \textbf{probabilistico}, quindi
lo si esegue molte volte: ogni esecuzione è uno \textbf{shot}. Mettendo
insieme i risultati di tutti gli shot si ottiene la distribuzione
statistica dei risultati. La \textbf{distribuzione degli shot} consiste
nel dividere gli shot di un circuito tra \textbf{più computer
quantistici} e poi combinare i dati per avere il quadro completo.

\end{boxdefinition}

Dato un circuito quantistico e dei requisiti specifici, il quantum
broker sceglie il miglior insieme di computer quantistici su cui
distribuire gli shot. Per ogni computer scelto individua i
\textbf{compilatori} più adatti e il \textbf{numero di shot} da
eseguire. I vantaggi sono:

\begin{itemize}
\tightlist
\item
  \textbf{Resilienza migliore}: dividere il lavoro tra sistemi
  quantistici diversi aumenta la tolleranza ai guasti.
\item
  \textbf{Alta personalizzazione}: si può controllare molto nel
  dettaglio come vengono gestiti i calcoli.
\item
  \textbf{Distribuzioni parziali}: ogni computer produce una
  distribuzione parziale, e le parziali si combinano per ottenere una
  visione completa dell'esecuzione del circuito, con la possibilità di
  metodi di unione più sofisticati.
\end{itemize}
